% ============================================================
%  ch18  附录 A §A.2 普遍局部无环性（ULA）
%  原文行 3538–3700
% ============================================================

\subsection{普遍局部无环性}\label{ULA sheaf}

设 $S$ 是概形，$s$ 是 $S$ 的几何点，以 $S_{(s)}$ 记 $S$ 在 $s$ 处的严格 Hensel 化。若 $t$ 是 $S_{(s)}$ 的几何点，则形式地写 $t\to s$，并称其为特化映射。

回忆 \cite{SGA4.5} 中的下述定义。

\begin{dfn}
设 $f:X\to S$ 是 $k$ 上有限型概形的态射。$X$ 上的 $\ell$-进复形 $\mF$ 称为关于 $f$ \textbf{局部无环}，若对每个几何点 $x\in X$ 与每个几何点 $t\in S_{(f(x))}$，自然映射 $R\Gamma(X_{(x)},\mF)\to R\Gamma( (X_{(x)})_t,\mF)$ 是同构，其中 $(X_{(x)})_t=X_{(x)}\times_{S_{(f(x))}}t$。若它在任意基变换 $S'\to S$ 之后仍局部无环，则称\textbf{普遍局部无环}（ULA）。
\end{dfn}

局部无环性可改写如下。设 $f:X\to S$ 如上，$\mF$ 是 $X$ 上的 $\ell$-进复形，$s$ 是 $S$ 的几何点，$t$ 是 $S_{(s)}$ 的几何点，$j_t: X\times_{S} t\to X$ 是自然映射。记
\[\Psi_{t\to s}(\mF):= ((j_t)_*j_t^*\mF)|_{X_s},\]
并把它看作沿特化映射 $t\to s$ 的邻近圈（回忆按约定 $(j_t)_*$ 是导出的）。事实上，若 $(S,s,\eta)$ 是严格局部 Hensel 链（即 $S$ 是局部严格 Hensel 离散赋值环的谱，$s$ 与 $\eta$ 是其闭点与一般点），$t$ 是 $\eta$ 上的几何点，则 $\Psi_{t\to s}(\mF)$ 按定义就是 $\mF$ 的一维邻近圈 $\Psi(\mF_\eta)$。

\begin{lem}\label{LA reform}
复形 $\mF$ 关于 $f$ 局部无环，当且仅当对如上每个 $s$ 与 $t$，自然映射
\begin{equation}\label{sp}
\mF|_{X_s}\to \Psi_{t\to s}(\mF)
\end{equation}
是同构。
\end{lem}

\begin{proof}
事实上，局部无环性说的正是上述映射在每个几何点 $x\in X_s$ 处诱导茎的同构。
\end{proof}

\begin{rmk}\label{LA explanation}
注意若 $(S,s,\eta)$ 是链，则引理 \ref{LA reform} 说 $\mF$ 关于 $f$ 局部无环当且仅当 $\mF_s\to \Psi(\mF_\eta)$ 是同构，等价地消没圈 $\Phi(\mF)$ 消没。因此对高维底 $S$，粗略地说，层 $\mF$ 关于 $f: X\to S$ 是 ULA，意味着沿 $S$ 的每个方向都没有消没圈。引入高维底上的邻近圈与消没圈函子即可使这一说法精确（参见 \cite[Example 1.7(b)]{Il3}）。
\end{rmk}

下面是所需的 ULA 层的基本性质。

\begin{thm}\label{ULA prop}~
\begin{enumerate}
\item 设 $S$ 是 $k$ 上的光滑簇。则 $S$ 上的复形 $\mF$ 关于恒等映射 ULA，当且仅当每个上同调层 $\mH^i(\mF)$ 是 $S$ 上的局部系。

\item 若 $f:X\to Y$ 是 $S$ 上的 proper 态射，且 $\mF$ 关于 $X\to S$ 是 ULA，则 $f_*\mF$ 关于 $S$ 是 ULA。

\item 若 $f:X\to Y$ 是 $S$ 上的光滑态射，则 $f^*\mF$ 关于 $X\to S$ 是 ULA，当且仅当 $\mF$ 关于 $Y\to S$ 是 ULA。

\item 设 $f_i:X_i\to S, i=1,2$ 是两个态射，$\mF_i$ 是 $X_i$ 上关于 $f_i$ 的 ULA 复形。则 $\mF_1\boxtimes_S\mF_2$ 关于 $f_1\times_Sf_2:X_1\times_SX_2\to S$ 是 ULA。

\item[(4)'] 设 $f_i:X_i\to S_i, i=1,2$，$\mF_i$ 是 $X_i$ 上关于 $f_i$ 的 ULA 复形。则 $\mF_1\boxtimes\mF_2$ 关于 $f_1\times f_2$ 是 ULA。
\end{enumerate}
\end{thm}

\begin{proof}
(1) 由定义可知 $S$ 上真正的层关于 $\id:S\to S$ 是 ULA 当且仅当它是局部系；一般情形对 $\mF$ 的上同调幅度归纳。

(2) 与 (3) 由函子 $\Psi_{t\to s}$ 与 proper 推前、光滑拉回交换推出。(4) 见 \cite[Corollary 2.5]{Il3}。(4)' 由 (4) 推出，因为对每个态射 $S\to S_1\times S_2$ 有
\[(X_1\times X_2)\times_{S_1\times S_2}S=(X_1\times_{S_1}S)\times_S(X_2\times_{S_2}S).\]
\end{proof}

下面讨论一个与 ULA 反常层有关的结果。

\begin{thm}\label{ULA perverse}
设 $f: X\to S$ 是 $k$-簇的态射，$S$ 光滑，$D\subset S$ 是光滑有效除子。设 $i: Z=f^{-1}(D)\to X$ 是闭嵌入，$j:U=X-Z\to X$ 是开补集。若 $X$ 上的复形 $\mF$ 关于 $f$ 是 ULA，且 $\mF|_U$ 反常，则
\[\mF\cong j_{!*}(\mF|_U), \quad i^!\mF[1]\cong i^*\mF[-1].\]
特别地 $\mF$ 也反常；此外 $i^*\mF[-1]$ 反常，且关于 $fi$ 是 ULA。
\end{thm}

\begin{proof}
问题在 \'etale 拓扑下是局部的，故可假设存在\emph{光滑}映射 $g: S\to \bA^1$ 使 $D=g^{-1}(0)$。由 \cite[Lemma 2.14]{SGA4.5}，$\mF$ 关于 $gf: X\to \bA^1$ 是 ULA。于是化归到 $f:X\to \bA^1$、$Z=f^{-1}(0)$ 的情形。由注 \ref{LA explanation}，$i^*\mF\to \Psi(j^*\mF)$ 是同构。由此
\[i^*\mF[-1]\cong \Psi(j^*\mF)[-1]=:\mC\]
反常（\cite[Corollary 4.5]{Il2}）。由 distinguished triangle
\[ j_!j^*\mF\to \mF\to i_*i^*\mF\to\]
可知 $\mF$ 只处于反常上同调次数 $-1$ 与 $0$。此外，由 distinguished triangle（见 \cite[Lemma 3.11]{SGA4.5} 的证明）
\[i^*j_*j^*\mF\to \Psi(j^*\mF)^P\stackrel{1-t}{\to} \Psi(j^*\mF)^P\to\]
可知反常上同调层 ${^p}\mH^r i^*j_*j^*\mF\cong \mC$（$r=-1,0$）且对 $r\neq -1,0$ 消没，其中 $P$ 记野惯性群，$t$ 记驯惯性群的生成元。于是由 distinguished triangle
\[i_*i^!\mF\to \mF\to j_*j^*\mF\to\]
得 $i^!\mF[1]\cong \mC$、$\mF$ 反常，且有反常层短正合列
\[0\to \mF\to j_*j^*\mF\to {^p}\mH^0j_*j^*\mF\to 0.\]
这又推出 $\mF\cong j_{!*}j^*\mF$。
\end{proof}

\begin{rmk}
当底 $S$ 光滑时，局部无环性还有 \cite[\S~5.1,~Appendix B]{BGa} 中引入的另一等价定义。按该定义，定理 \ref{ULA prop} (iv) 由 Gaitsgory 证明（私人通信），定理 \ref{ULA prop} 其余部分与定理 \ref{ULA perverse} 由 Reich 在 \cite{Rei} 中证明。
\end{rmk}
